\section{Results}
\label{sec:results}

We present the nDCG@10 scores on the \textit{dev} split for our three selected domains. Table \ref{tab:results} consolidates the results across all baselines.

\begin{table}[h]
\centering
\begin{tabular}{lcccc}
\toprule
\textbf{Method} & \textbf{quant} & \textbf{law} & \textbf{workplace} & \textbf{Macro} \\
\midrule
BM25 (k1=0.9, b=0.4) & 0.0237 & 0.0574 & 0.0230 & 0.0347 \\
TF-IDF (Cosine) & 0.0053 & 0.0463 & 0.0094 & 0.0203 \\
Dense (bge-small frozen) & 0.0919 & 0.3697 & 0.2676 & 0.2431 \\
LightGBM Reranker (Top-50) & 0.0305 & 0.1340 & 0.0167 & 0.0604 \\
Hybrid (RRF BM25+Dense) & 0.0022 & 0.0467 & 0.0124 & 0.0204 \\
\bottomrule
\end{tabular}
\caption{Consolidated nDCG@10 Results on Dev Split. \textit{Note: scores may update upon final pipeline execution.}}
\label{tab:results}
\end{table}

The LightGBM pointwise reranker outperforms the BM25 baseline on most domains due to its explicit temporal matching features. Surprisingly, the frozen BGE-small dense retriever achieved the highest overall score (0.2431 Macro nDCG@10). This indicates that bridging the vocabulary gap via semantic representations is highly beneficial for these specific domains, compensating for its lack of explicit temporal logic. Hybrid RRF fusion failed to properly combine the signals due to scale mismatches between BM25 scores and Dense similarities without calibration.
